\documentclass[12pt]{article}
\usepackage[utf8]{inputenc}
\usepackage[T1]{fontenc}
\usepackage{hyperref}
\title{Pashai, Southwest Language Resources for K-12 Curriculum Development}
\author{Compiled by Claude Sonnet 5.5 (claude-sonnet-5-5) for Andr\'as Kornai}
\date{}
\begin{document}
\maketitle

\section{Introduction}
Southwest Pashai (ISO 639-3 psh) is an Indo-Aryan language of about 100,000 speakers in the Hindu Kush valleys northeast of Kabul (Kapisa, Tagab, Laghman). Dialects named include Ishpi and Tagau. It is not known to be taught in schools, and the searches found no school materials. Do not confuse it with Southeastern Pashayi (psi) of Dara-i-Nur.

\section{Primary Educational Resources}
None found. Literacy is very low (for Southeast Pashayi, under one percent in the first language, per Joshua Project).

\section{Linguistic and Reference Materials}
\begin{itemize}
\item Georg Morgenstierne, \emph{The Pashai Language} (texts, dictionary, grammar; 1944--1967, 2nd ed.\ 1973), the most thorough published study of Pashayi overall. Bibliographic and review pages: \url{https://fid4sa-repository.ub.uni-heidelberg.de/3173} and \url{https://www.iranicaonline.org/articles/morgenstierne-georg-valentin-von-munthe-af/}
\item Max Planck numeral database page for Southwest Pashai (vigesimal counting; Laghman dialect). \url{https://lingweb.eva.mpg.de/channumerals/Pashai-Southwest.htm}
\item Illustration of the IPA, Southeastern Pashayi (a different variety, psi), Journal of the International Phonetic Association. \url{https://www.cambridge.org/core/product/8CDDFBEC9998A428ACA1FFB023555626/core-reader}
\item Wikipedia, Pashayi languages. \url{https://en.wikipedia.org/wiki/Pashayi_languages}
\end{itemize}

\section{Religious and Community Resources}
\begin{itemize}
\item Joshua Project, Southwest Pashayi in Afghanistan (Scripture status: translation started). \url{https://www.joshuaproject.net/people_groups/15072/AF}
\item Joshua Project language page psh. \url{https://www.joshuaproject.net/index.php/languages/psh}
\end{itemize}

\section{Digital and Audio Resources}
\begin{itemize}
\item Global Recordings Network, Pashai, Southwest language page (no recordings in this variety when checked). \url{https://globalrecordings.net/en/language/psh}
\item Global Recordings Network, Pashai: Kohistani \& Shahi (related varieties). \url{https://www.globalrecordings.net/en/language/3317}
\end{itemize}

\section{Curriculum Development Context}
No orthography, primers or school programmes for Southwest Pashai were found. A curriculum would start from orthography development and literacy work; Morgenstierne's materials and a Pashai Language Development Project (mentioned in Wikipedia, not independently verified) are possible starting points.
\end{document}